\chapter{WRF-Fire}\label{ch:fire}

WRF-Fire \citep{coen2013,mandel2011} moves a fire front over the landscape at the rate of spread of
\cref{ch:firebasics}, burns the fuel behind it, and feeds the heat and moisture into the atmosphere. In turn the
atmosphere changes the wind that drives the front. This chapter describes its numerics as in the WRF~4.6.0 source:
\begin{itemize}
  \item \code{phys/module_fr_fire_driver.F} and \code{module_fr_fire_driver_wrf.F}: the drivers;
  \item \code{module_fr_fire_model.F}: the sequence of a fire step;
  \item \code{module_fr_fire_core.F}: the level set, ignition and fuel;
  \item \code{module_fr_fire_phys.F}: rate of spread and fluxes;
  \item \code{module_fr_fire_atm.F}: the coupling to the atmosphere;
  \item \code{module_fr_fire_util.F}: utilities.
\end{itemize}

\section{The fire grid}

The fire lives on a grid finer than the atmosphere's by \code{sr_x}$\times$\code{sr_y} (4$\times$4 in the case), so
each 100~m atmospheric cell holds 16 fire cells of 25~m. Fire variables are stored at the fire-grid nodes, with
memory dimensions that include a halo. The domain's outermost fire rows and columns are \emph{ghost} cells: the
fuel calculation reads them, so their values are part of the state (\cref{ch:subsys}). The case's d02 has
$3244\times3244$ fire nodes, so a single fire field holds 10.5 million values. Several such fields are updated every
d02 step.

\section{From the atmosphere to the fire}

The fire spreads with the wind near the ground, not at the first model level (about 10--20~m up). Before each fire
step, \src{phys/module_fr_fire_driver.F}{interpolate_atm2fire} computes it:
\begin{enumerate}
  \item The heights of the model levels at $u$ and $v$ points are computed.
  \item The wind at the fire wind height (\code{fire_wind_height}, 1~m in the case) is interpolated from the two
    levels that bracket it. With \code{fire_lsm_zcoupling}, a logarithmic profile with the surface roughness
    $z_0$ is used:
    \begin{equation}
      u(z) = u(z_k)\,\frac{\ln(z/z_0)}{\ln(z_k/z_0)} .
    \end{equation}
    The search for the bracketing level stops at the first level above the target height (an early exit, written with
    \code{GOTO} in the source).
  \item The values are extended to the domain edges (\src{phys/module_fr_fire_util.F}{continue_at_boundary}): first
    along $j$ strips, then $i$ strips, then the corners.
  \item The winds are interpolated bilinearly to the fire nodes (\src{phys/module_fr_fire_util.F}{interpolate_2d}).
    Each coarse cell writes its $4\times4$ fine nodes.
\end{enumerate}

\section{The level-set equation}

The fire front is the zero contour of a level-set function $\psi(x, y, t)$ (\code{lfn}): burning or burned where
$\psi \le 0$, unburned where $\psi > 0$. The front moves in its normal direction at the rate of spread $R$, so
\begin{equation}
  \pd{\psi}{t} + R(x, y, \mathbf{n}, t)\,|\nabla\psi| = 0, \qquad \mathbf{n} = \frac{\nabla\psi}{|\nabla\psi|} .
  \label{eq:levelset}
\end{equation}
$R$ depends on the direction through the wind and slope components normal to the front (\cref{alg:ros}).
\src{phys/module_fr_fire_core.F}{tend_ls} computes the right-hand side $-R|\nabla\psi|$ at every node.

\paragraph{Discretizing $|\nabla\psi|$.} The gradient must be upwinded: information flows from the burned side to the
unburned side. WRF-Fire offers nine schemes (\code{fire_upwinding}). The case uses option 9, a hybrid:
\begin{itemize}
  \item a fifth-order WENO reconstruction of the one-sided derivatives \citep{jiang1996} within a band of
    \code{fire_lsm_band_ngp} nodes (4) around the front, where accuracy matters (\code{select_weno5});
  \item first-order ENO elsewhere and within 10 nodes of the domain edge (\code{select_eno}).
\end{itemize}
The gradient magnitude is
$|\nabla\psi| = \sqrt{(\delta_x\psi)^2 + (\delta_y\psi)^2}$ with the upwinded one-sided differences.

\paragraph{Artificial viscosity.} To damp oscillations of the front, \code{tend_ls} adds a term proportional to
\code{fire_viscosity} (0.4 in the case), the grid spacing and the second differences of $\psi$, within
\code{fire_viscosity_ngp} nodes of the front.

\paragraph{Time step bound.} \code{tend_ls} also returns the largest $R|\nabla\psi|/\Delta x$, a reduction over the
grid. From it the driver checks that the fire time step is stable.

\section{Time stepping}

The level set is advanced with the same third-order Runge--Kutta scheme as the atmosphere
(\src{phys/module_fr_fire_core.F}{prop_ls_rk3}):
\begin{equation}
  \psi_1 = \psi^n + \tfrac{\Delta t}{3}\,T(\psi^n), \qquad
  \psi_2 = \psi^n + \tfrac{\Delta t}{2}\,T(\psi_1), \qquad
  \psi^{n+1} = \psi^n + \Delta t\,T(\psi_2),
\end{equation}
where $T$ is the tendency of \cref{eq:levelset}. Each evaluation is preceded by \code{continue_at_boundary}. With
\code{fire_grows_only = 1}, the new level set is $\min(\psi^{n+1}, \psi^n)$, so a burned area never shrinks.

\paragraph{Ignition time.} Where $\psi$ changes sign during the step, the node ignites. Its ignition time is
interpolated linearly in time (\src{phys/module_fr_fire_core.F}{tign_update}):
\begin{equation}
  t_i = t^n + \Delta t\,\frac{\psi^n}{\psi^n - \psi^{n+1}} .
\end{equation}
The ignition time \code{tign_g} is the main output of a fire simulation. Its contours are the fire perimeters.

\section{Reinitialization}

The level-set function drifts away from a signed distance function, and then its gradient becomes too steep or too
flat near the front. With \code{fire_lsm_reinit}, WRF-Fire restores the distance property every step by solving, in
pseudo-time $\tau$, the reinitialization equation of \citet{sussman1994}:
\begin{equation}
  \pd{\psi}{\tau} + S(\psi_0)\,\bigl(|\nabla\psi| - 1\bigr) = 0, \qquad
  S(\psi_0) = \frac{\psi_0}{\sqrt{\psi_0^2 + \Delta x^2}} .
\end{equation}
The smoothed sign $S$ keeps the zero contour in place. The case does one RK3 iteration of it
(\code{fire_lsm_reinit_iter = 1}) with hybrid WENO5/ENO1 differences
(\src{phys/module_fr_fire_core.F}{reinit_ls_rk3}, \code{advance_ls_reinit}). The result is again limited by
$\min(\cdot, \psi)$, so reinitialization cannot extinguish the fire.

\section{Ignition}

Ignitions are given in the namelist: position, radius, start and end time, and the rate at which the ignition circle
grows. \src{phys/module_fr_fire_core.F}{ignite_fire} computes, for every node, the distance $d$ to the ignition point
and a level set of the growing circle:
\begin{equation}
  \psi_{\mathrm{ign}} = d - r(t),
\end{equation}
and sets $\psi \gets \min(\psi, \psi_{\mathrm{ign}})$. Newly ignited nodes get their ignition time from the growth
rate. The case ignites one circle of 100~m radius between 8280~s and 9080~s of model time.

\section{Fuel consumption}

\src{phys/module_fr_fire_core.F}{fuel_left} computes the fraction of the original fuel left in each fire cell. With
\code{fire_fuel_left_method = 1}, each cell is divided into $2\times2$ sub-cells. In each sub-cell
\code{fuel_left_cell_1} estimates:
\begin{itemize}
  \item the burned area fraction, from the level set at the sub-cell corners;
  \item the average of $\exp(-(t - t_i)/T_f)$ over that area, from the ignition times at the corners.
\end{itemize}
The cell's value is the average over its four sub-cells, taken in the source's order. The difference between the
fuel fraction at the start and the end of the step is the fuel burned during the step.

\section{Heat release and coupling}

\src{phys/module_fr_fire_phys.F}{heat_fluxes} turns the burned fuel into fluxes at the ground. With the dry mass burned
$\Delta m = w_0 \times \text{(fraction burned in the step)}$ and the fuel moisture fraction
$b = M_f/(1 + M_f)$:
\begin{equation}
  H_g = \frac{\Delta m}{\Delta t}\,(1 - b)\,h_c, \qquad
  E_g = \bigl(b + 0.56\,(1 - b)\bigr)\,\frac{\Delta m}{\Delta t}\,L_v ,
\end{equation}
where $h_c$ is the heat of combustion (\code{cmbcnst}) and 0.56 the water produced per unit dry fuel burned.
\code{sum_2d_cells} sums the fire-grid fluxes onto the atmospheric cells.
\src{phys/module_fr_fire_atm.F}{fire_tendency} then distributes them in the vertical with an exponential decay over
the extinction depth $\alpha_g$ (\code{fire_ext_grnd}):
\begin{equation}
  \left.\pd{\theta}{t}\right|_{k} \propto \frac{H_g}{\rho\,c_p}\,\frac{e^{-z_{k}/\alpha_g} - e^{-z_{k+1}/\alpha_g}}{\Delta z_k},
\end{equation}
and the same for moisture. These tendencies (\code{rthfrten}, \code{rqvfrten}) are added to the atmosphere's like any
physics tendency (\cref{ch:phys}). The heat drives the fire's convective column, which draws air toward the front:
the fire-atmosphere feedback.

\begin{algorithm}[ht]
\caption{One fire step on d02 (inside RK stage 1, \code{fire_driver_em_step} $\to$ \code{fire_model}).}
\label{alg:fire-step}
\begin{algorithmic}[1]
\State wind at the fire wind height on the fire grid (\code{interpolate_atm2fire})
\State ignition (\code{ignite_fire}) if within an ignition period
\State level set: three RK stages of \code{tend_ls} with boundary continuation (\code{prop_ls_rk3})
\State ignition times of newly burning nodes (\code{tign_update}); boundary guard; flame length
\State reinitialization (\code{reinit_ls_rk3}), limited by the old level set
\State fuel left and fuel burned (\code{fuel_left})
\State heat and moisture fluxes (\code{heat_fluxes}); sums onto atmospheric cells (\code{sum_2d_cells})
\State tendencies of $\theta$ and $q_v$ in the atmospheric columns (\code{fire_tendency})
\end{algorithmic}
\end{algorithm}

\begin{portnote}
The fire step is a sequence of two-dimensional stencils on a large grid, ideal for a GPU, with four complications.
\begin{itemize}
  \item \textbf{Branches.} The WENO/ENO choice, the ignition test and the fuel test depend on the position relative
    to the front. Neighbouring threads take different branches (divergence, \cref{ch:warp}).
  \item \textbf{Ghost cells.} The fuel calculation reads the ghost rows and columns. On the CPU they keep their values
    from allocation, so the device copy must hold the same values (test T-FIRE-GHOST).
  \item \textbf{Pointer components.} The fuel parameters are pointer components of a derived type (\texttt{fp\%r\_0},
    \dots). The device needs that type copied, with its pointers attached to device memory.
  \item \textbf{Sensitivity.} Bitwise equality of the level set is the whole point of the port. A single bit in
    $\psi$ near zero decides whether a node ignites in this step or the next. The fire windows T-FIRE-IGN and
    T-FIRE-WIN compare every fire array at every step, and the burned cells at each history time
    (\cref{ch:subsys}).
\end{itemize}
\end{portnote}
